\chapter{Proposed Methodology}
\label{ch:methodology}

This chapter provides a comprehensive, component-level breakdown of the automated vessel segmentation and morphological analysis pipeline developed in this thesis. The methodology is broadly partitioned into three primary stages: dataset preprocessing and augmentation, semantic segmentation using advanced deep learning architectures (UNet++ and DeepLabV3+), and post-segmentation mathematical skeletonization for topological feature extraction.

\section{Dataset Description and Preprocessing}
The foundation of any deep learning model lies in the quality, diversity, and statistical distribution of its training data. The primary dataset utilized in this research consists of 637 high-resolution RGB angiography images, partitioned into a training set of 500 images and an independent, hold-out testing set of 137 images. 

Each RGB image, denoted as $I \in \mathbb{R}^{H \times W \times 3}$, is accompanied by a binary ground truth mask $G \in \{0,1\}^{H \times W}$ generated through automated pre-processing. The preparation of these ground truth masks followed a systematic automated pipeline. First, the green color channel was extracted from the raw RGB images to maximize the contrast of the hemoglobin-rich vasculature against the surrounding chorioallantoic membrane tissue. An automated binarization was then performed using localized adaptive Gaussian thresholding to segment the vessel boundaries. Because manual pixel-wise refinement of 637 high-resolution images containing dense vascular structures is practically infeasible within experimental constraints, these automated thresholded masks were utilized directly as the training targets. While this automated step introduced local background artifacts (salt-and-pepper noise) in the ground truth, utilizing these weakly supervised targets allowed for evaluating the model's intrinsic capacity to generalize and filter out high-frequency noise during training. Finally, to ensure computational efficiency and standardized batch processing during network training, all input images and their corresponding masks were resized to a uniform spatial resolution of $256 \times 256$ pixels using bilinear interpolation for the images and nearest-neighbor interpolation for the masks to preserve discrete binary boundaries.

Following spatial normalization, the pixel intensities of the RGB images are normalized from the range $[0, 255]$ to the range $[0.0, 1.0]$. This vital preprocessing step stabilizes the gradients during backpropagation and accelerates the convergence of the neural networks by ensuring that the input features operate on a similar numerical scale.

\section{Architecture 1: The UNet++ Framework}
The first architecture explored in this thesis is UNet++, specifically designed to overcome the semantic gap inherent in the skip connections of standard U-Net models. We implement and evaluate two distinct variants of UNet++: a pure, from-scratch 4-level encoder, and a hybrid architecture utilizing a frozen, pre-trained ResNet50 backbone.

\subsection{Dense Skip Pathways Formulation}
The core innovation of UNet++ lies in its dense, nested skip pathways. Let $x^{i,j}$ denote the output feature map of node $X_{i,j}$ in the UNet++ grid, where $i \in \{0, 1, 2, 3, 4\}$ indexes the down-sampling layer along the encoder (depth), and $j$ indexes the convolution layer along the dense skip pathway.

The feature map $x^{i,j}$ is mathematically computed as follows:
\begin{equation}
    x^{i,j} = \begin{cases}
        \mathcal{H} \left( \mathcal{D}(x^{i-1, j}) \right) & \text{if } j = 0 \\
        \mathcal{H} \left( \left[ \left[ x^{i,k} \right]_{k=0}^{j-1}, \mathcal{U}(x^{i+1, j-1}) \right] \right) & \text{if } j > 0
    \end{cases}
\end{equation}
where $\mathcal{H}(\cdot)$ represents a convolution block consisting of two consecutive sequences of $3 \times 3$ Convolution $\rightarrow$ Batch Normalization $\rightarrow$ ReLU activation, followed by an optional spatial dropout layer. $\mathcal{D}(\cdot)$ denotes a $2 \times 2$ Max Pooling operation for spatial down-sampling (stride 2), and $\mathcal{U}(\cdot)$ denotes a $2 \times 2$ bilinear Up-Sampling operation. The operation $[\cdot]$ denotes the channel-wise concatenation of feature maps.

For instance, node $X_{0,3}$ receives the concatenated outputs of all its same-level predecessors ($X_{0,0}, X_{0,1}, X_{0,2}$) as well as the up-sampled output of the deeper node $X_{1,2}$. This dense concatenation ensures that by the time the feature maps reach the final full-resolution decoder node $X_{0,4}$, the network has fused features of highly varying semantic depth, thereby preserving the spatial continuity of thin vessel segments.

\subsection{UNet++ with ResNet50 Encoder}
To leverage transfer learning, we graft the dense UNet++ decoder onto an ImageNet-pretrained ResNet50 backbone. The ResNet50 encoder processes the $256 \times 256$ input image through five primary stages, generating multi-scale feature maps at strides of 2, 4, 8, 16, and 32. To prevent catastrophic forgetting and over-parameterization on our relatively small dataset of 500 training images, the weights of the early layers of the ResNet50 backbone (specifically the first 80 layers) are frozen during training.

The output feature maps from the ResNet50 blocks (e.g., \texttt{conv1\_relu}, \texttt{conv2\_block3\_out}, etc.) are projected via $3 \times 3$ convolutions to a uniform channel depth (128 channels). These projected multi-scale features serve as the $j=0$ column of the UNet++ grid, seamlessly feeding into the dense skip pathways described previously.

\begin{figure}[H]
    \centering
    \begin{tikzpicture}[
        node distance=1.0cm and 1.0cm,
        enc/.style={rectangle, draw=blue!70, fill=blue!15, minimum width=1.0cm, minimum height=0.7cm, align=center, font=\footnotesize, rounded corners=2pt},
        dec/.style={rectangle, draw=green!70!black, fill=green!15, minimum width=1.0cm, minimum height=0.7cm, align=center, font=\footnotesize, rounded corners=2pt},
        skip/.style={rectangle, draw=orange!70, fill=orange!15, minimum width=1.0cm, minimum height=0.7cm, align=center, font=\footnotesize, rounded corners=2pt},
        arr/.style={-{Stealth[length=2mm]}, thick},
        darr/.style={-{Stealth[length=2mm]}, thick, dashed}
    ]
        % Encoder column (j=0)
        \node[enc] (x00) {$X^{0,0}$};
        \node[enc, below=of x00] (x10) {$X^{1,0}$};
        \node[enc, below=of x10] (x20) {$X^{2,0}$};
        \node[enc, below=of x20] (x30) {$X^{3,0}$};
        \node[enc, below=of x30] (x40) {$X^{4,0}$};

        % Dense skip nodes
        \node[skip, right=of x00] (x01) {$X^{0,1}$};
        \node[skip, right=of x10] (x11) {$X^{1,1}$};
        \node[skip, right=of x20] (x21) {$X^{2,1}$};
        \node[skip, right=of x30] (x31) {$X^{3,1}$};

        \node[skip, right=of x01] (x02) {$X^{0,2}$};
        \node[skip, right=of x11] (x12) {$X^{1,2}$};
        \node[skip, right=of x21] (x22) {$X^{2,2}$};

        \node[skip, right=of x02] (x03) {$X^{0,3}$};
        \node[skip, right=of x12] (x13) {$X^{1,3}$};

        \node[dec, right=of x03] (x04) {$X^{0,4}$};

        % Encoder downsampling arrows
        \draw[arr, blue!60] (x00) -- (x10);
        \draw[arr, blue!60] (x10) -- (x20);
        \draw[arr, blue!60] (x20) -- (x30);
        \draw[arr, blue!60] (x30) -- (x40);

        % Horizontal skip connections
        \draw[arr] (x00) -- (x01);
        \draw[arr] (x10) -- (x11);
        \draw[arr] (x20) -- (x21);
        \draw[arr] (x30) -- (x31);
        \draw[arr] (x01) -- (x02);
        \draw[arr] (x11) -- (x12);
        \draw[arr] (x21) -- (x22);
        \draw[arr] (x02) -- (x03);
        \draw[arr] (x12) -- (x13);
        \draw[arr] (x03) -- (x04);

        % Upsampling diagonal arrows
        \draw[arr, red!60] (x10) -- (x01);
        \draw[arr, red!60] (x20) -- (x11);
        \draw[arr, red!60] (x30) -- (x21);
        \draw[arr, red!60] (x40) -- (x31);
        \draw[arr, red!60] (x11) -- (x02);
        \draw[arr, red!60] (x21) -- (x12);
        \draw[arr, red!60] (x31) -- (x22);
        \draw[arr, red!60] (x12) -- (x03);
        \draw[arr, red!60] (x22) -- (x13);
        \draw[arr, red!60] (x13) -- (x04);
    \end{tikzpicture}
    \caption{Architectural diagram of the UNet++ dense skip pathway topology. Blue nodes represent the encoder column ($j=0$), orange nodes represent intermediate dense skip nodes, and the green node represents the final decoder output ($X^{0,4}$). Blue arrows denote max-pooling downsampling, black arrows denote horizontal dense concatenation, and red arrows denote bilinear upsampling paths.}
    \label{fig:unetpp_architecture}
\end{figure}

\section{Architecture 2: The DeepLabV3+ Framework}
The second architecture implemented is DeepLabV3+, which approaches the multi-scale challenge not through dense decoder pathways, but through advanced atrous (dilated) convolutions within the encoder bottleneck. 

\subsection{Atrous Spatial Pyramid Pooling (ASPP)}
The core component of our DeepLabV3+ implementation is the Atrous Spatial Pyramid Pooling (ASPP) module. Let $y = f(x)$ denote a 2D discrete convolution. For a standard convolution, the output $y[i,j]$ over a feature map $x$ with a filter $w$ of size $K \times K$ is:
\begin{equation}
    y[i,j] = \sum_{m,n}^{K} x[i+m, j+n] \cdot w[m,n]
\end{equation}
An atrous convolution introduces a dilation rate $r$, modifying the operation to:
\begin{equation}
    y[i,j] = \sum_{m,n}^{K} x[i+r \cdot m, j+r \cdot n] \cdot w[m,n]
\end{equation}
When $r=1$, this reverts to a standard convolution. When $r>1$, the filter's receptive field is expanded without increasing the number of trainable parameters or down-sampling the spatial resolution.

In our methodology, the ResNet50 backbone processes the image down to a stride of 16 (output shape $16 \times 16$). The ASPP module then processes this $16 \times 16$ feature map simultaneously through five parallel branches:
\begin{enumerate}
    \item A $1 \times 1$ convolution branch.
    \item A $3 \times 3$ atrous convolution with dilation rate $r = 6$.
    \item A $3 \times 3$ atrous convolution with dilation rate $r = 12$.
    \item A $3 \times 3$ atrous convolution with dilation rate $r = 18$.
    \item An Image Pooling branch, which applies Global Average Pooling over the entire $16 \times 16$ spatial dimension, followed by a $1 \times 1$ convolution, and is then bilinearly up-sampled back to $16 \times 16$.
\end{enumerate}
The outputs of all five branches are concatenated channel-wise and projected through a final $1 \times 1$ convolution block. This ASPP tensor encapsulates rich semantic context across highly varying spatial scales.

\subsection{DeepLabV3+ Decoder Design}
The ASPP output tensor is initially up-sampled by a factor of 4 (via bilinear interpolation) to a spatial resolution of $64 \times 64$. Concurrently, low-level spatial features from the ResNet50 encoder (specifically the output from \texttt{conv2\_block3\_out}, which also sits at a $64 \times 64$ resolution) are projected down to 48 channels via a $1 \times 1$ convolution to reduce their dimensional dominance.

These low-level spatial features are then concatenated with the up-sampled ASPP features. This fused tensor undergoes a series of two $3 \times 3$ convolutions to refine object boundaries. Finally, the tensor is up-sampled by another factor of 4 to match the original $256 \times 256$ input resolution. A final $1 \times 1$ convolution with a Sigmoid activation function maps the dense tensor to a binary probability map $P \in [0,1]^{256 \times 256}$.

\begin{figure}[H]
    \centering
    \begin{tikzpicture}[
        node distance=0.5cm and 0.8cm,
        block/.style={rectangle, draw=black!60, fill=#1, text width=2.4cm, minimum height=0.9cm, align=center, rounded corners=2pt, font=\footnotesize},
        arr/.style={-{Stealth[length=2.5mm]}, thick, draw=black!50}
    ]
        \node[block=blue!12] (input) {Input Image\\$256 \times 256$};
        \node[block=blue!20, right=of input] (backbone) {ResNet50\\Backbone};
        \node[block=red!15, right=of backbone] (aspp) {ASPP Module\\$(r=6,12,18)$};
        \node[block=green!15, below=1.2cm of backbone] (lowlevel) {Low-Level\\Features ($64^2$)};
        \node[block=orange!15, right=of aspp] (concat) {Concatenate\\\& Refine};
        \node[block=purple!12, right=of concat] (output) {Output Mask\\$256 \times 256$};

        \draw[arr] (input) -- (backbone);
        \draw[arr] (backbone) -- (aspp);
        \draw[arr] (aspp) -- (concat);
        \draw[arr] (backbone) |- (lowlevel);
        \draw[arr] (lowlevel) -| (concat);
        \draw[arr] (concat) -- (output);
    \end{tikzpicture}
    \caption{Architectural diagram of the DeepLabV3+ framework. The ResNet50 encoder extracts multi-scale features, which are processed by the ASPP module with parallel atrous convolutions at dilation rates of 6, 12, and 18. The decoder fuses the ASPP output with low-level encoder features to refine spatial boundaries before producing the final segmentation mask.}
    \label{fig:deeplabv3p_architecture}
\end{figure}

\section{Post-Processing: Topological Skeletonization}
The primary objective of automated angiography analysis is not merely to generate a visual mask, but to extract actionable clinical numbers. To transition from pixel-level predictions to morphological metrics, we designed a downstream post-processing pipeline centered on mathematical skeletonization.

\subsection{Thresholding and Binarization}
The continuous probability map $P$ generated by either UNet++ or DeepLabV3+ is first thresholded into a discrete binary mask $M$:
\begin{equation}
    M(x,y) = \begin{cases}
        1 & \text{if } P(x,y) \ge 0.5 \\
        0 & \text{if } P(x,y) < 0.5
    \end{cases}
\end{equation}

\subsection{Morphological Thinning}
The binary mask $M$ is then passed through an iterative thinning algorithm (such as the classical Zhang-Suen algorithm). Let $\mathcal{S}(M)$ denote the skeletonization operation. The algorithm iteratively erodes the boundary pixels of the foreground vessel structures until the structures are exactly one pixel wide, resulting in the topological skeleton graph $S = \mathcal{S}(M)$. 

This skeleton preserves the fundamental Euler characteristic of the vascular tree. Because the width of the vessels has been collapsed to a single pixel line, we can now reliably count distinct vascular structures by identifying connected components within the graph $S$. 

\subsection{Feature Extraction: Junctions and Networks}
To extract morphological features, we traverse the skeleton graph $S$ using a $3 \times 3$ local neighborhood operation. For every foreground pixel $p_i \in S$, we count the number of adjacent foreground pixels in its 8-connected neighborhood $N_8(p_i)$.
\begin{itemize}
    \item \textbf{End Points:} If $\sum N_8(p_i) == 1$, the pixel is classified as a terminal end point of a capillary.
    \item \textbf{Branch Points (Bifurcations):} If $\sum N_8(p_i) \ge 3$, the pixel is classified as a bifurcation or crossover junction.
\end{itemize}

By removing the branch points from the skeleton graph $S$, the network decomposes into a set of disjoint, isolated vessel segments. Conversely, by applying a standard Connected Component Labeling (CCL) algorithm to the full skeleton $S$, we can accurately count the number of \textit{distinct, non-contiguous vessel networks} present in the angiogram. These quantitative features (Network Count, Branch Point Count, End Point Count) are automatically logged to a CSV report (e.g., \texttt{vessel\_counts.csv}), providing clinicians with a robust, objective numerical analysis of the patient's vascular health.

\section{Training Algorithm}
The complete training procedure for both architectures is formalized in Algorithm~\ref{alg:training}.

\begin{algorithm}[H]
\SetAlgoLined
\KwIn{Training set $\mathcal{D}_{train} = \{(I_k, G_k)\}_{k=1}^{500}$, Validation set $\mathcal{D}_{val}$}
\KwOut{Trained model weights $\theta^*$}
\textbf{Initialize:} Model $\mathcal{F}_\theta$, Adam optimizer ($\alpha_0 = 10^{-4}$), patience $p=10$\;
\For{epoch $= 1$ \KwTo $50$}{
    \ForEach{mini-batch $(I_b, G_b)$ from $\mathcal{D}_{train}$}{
        $(I'_b, G'_b) \leftarrow \texttt{Augment}(I_b, G_b)$\;
        $P_b \leftarrow \mathcal{F}_\theta(I'_b)$\;
        $\mathcal{L} \leftarrow \mathcal{L}_{BCE+Dice}(P_b, G'_b)$\;
        $\theta \leftarrow \theta - \alpha \nabla_\theta \mathcal{L}$\;
    }
    Compute $\mathcal{L}_{val}$ on $\mathcal{D}_{val}$\;
    \If{$\mathcal{L}_{val}$ improved}{Save $\theta$ as best weights\;}
    \If{no improvement for $p$ epochs}{\textbf{break} (Early Stopping)\;}
    \If{no improvement for $5$ epochs}{$\alpha \leftarrow 0.2 \cdot \alpha$ (ReduceLROnPlateau)\;}
}
Restore best weights $\theta^*$\;
\caption{Training Procedure for Vessel Segmentation Models}
\label{alg:training}
\end{algorithm}

\section{Skeletonization and Feature Extraction Algorithm}
The post-processing morphological analysis pipeline is formalized in Algorithm~\ref{alg:skeleton}.

\begin{algorithm}[H]
\SetAlgoLined
\KwIn{Predicted probability map $P \in [0,1]^{H \times W}$}
\KwOut{Network count $N_c$, Branch points $B_p$, End points $E_p$}
$M \leftarrow \texttt{Threshold}(P, \tau=0.5)$\;
$S \leftarrow \texttt{ZhangSuenThinning}(M)$ \tcp*{Reduce to 1-pixel skeleton}
$B_p \leftarrow 0$, $E_p \leftarrow 0$\;
\ForEach{foreground pixel $p_i \in S$}{
    $n \leftarrow |N_8(p_i)|$ \tcp*{Count 8-connected neighbors}
    \If{$n = 1$}{$E_p \leftarrow E_p + 1$ \tcp*{End point}}
    \If{$n \geq 3$}{$B_p \leftarrow B_p + 1$ \tcp*{Branch point}}
}
$N_c \leftarrow \texttt{ConnectedComponents}(S)$\;
\Return{$N_c, B_p, E_p$}\;
\caption{Morphological Skeletonization and Feature Extraction}
\label{alg:skeleton}
\end{algorithm}
